
\subsection{Symmetric and Antisymmetric Sub-Manifolds}
\label{subsec:introduce-symmetric-antisymmetric-branches-and-saddle-point}
This subsection introduces the symmetric and antisymmetric coordinates that separate the teacher branch from the antipodal saddle branch.

We first collect the active odd moments into a single finite-dimensional tensor space.

Define the active odd raw-moment space
\begin{equation}
    \mathcal H_{\mathcal O}
    \coloneq
    \mathcal H_{\mathcal O}^{\mathrm{raw}}
    =
    \bigoplus_{n\in\mathcal O}\Sym^n(\R^{d_w}),
    \qquad
    \|\mathbf X\|_{\mathcal H_{\mathcal O}}^2
    \coloneq
    \|\mathbf X\|_{\mathrm{raw}}^2
    =
    \sum_{n\in\mathcal O}\|X_n\|_{\HS}^2.
\end{equation}

The weighted norm records the size of active moments as seen through the
activation coefficients:
\begin{equation}
\label{eq:def-weighted-active-norm}
    \|\mathbf X\|_{a}^2
    \coloneq
    \sum_{n\in\mathcal O}\|a_n\|_2^2\|X_n\|_{\HS}^2.
\end{equation}

We now pass from the two raw moment sequences to symmetric and antisymmetric coordinates.

\begin{definition}
Introduce the two active student blocks and the teacher block
\begin{equation}
    \mathbf T_i
    \coloneq
    (T_{i,n})_{n\in \mathcal O}
    \in \mathcal H_{\mathcal O},\quad i=1,2,
    \quad 
    \mathbf T^o
    \coloneq
    (T_n^o)_{n\in \mathcal O}
    \in \mathcal H_{\mathcal O},
\end{equation}
where $T_{i,n}=\int w^{\otimes n}\,d\mu_i(w)$ and
$T_n^o=\int w^{\otimes n}\,d\nu^o(w)$ as in
\cref{assumption:odd-symmetric-teacher}. Recall the symmetric and
antisymmetric blocks
\begin{equation}
\label{eq:mf-symmetric-antisymmetric-active-blocks}
    \mathbf S
    \coloneq
    (S_n)_{n\in \mathcal O},
    \qquad
    \mathbf A
    \coloneq
    (A_n)_{n\in \mathcal O}.
\end{equation}
with $S_n=(T_{1,n}+T_{2,n})/2$ and
$A_n=(T_{1,n}-T_{2,n})/2$. Thus
\begin{equation}
    T_{1,n}=S_n+A_n,
    \qquad
    T_{2,n}=S_n-A_n,
    \qquad
    n\in \mathcal O.
\end{equation}

The finite active moment map used in the branch geometry is
\begin{equation}
\label{eq:mf-finite-active-moment-map}
    \Theta(\mu_1,\mu_2)
    \coloneq
    \bigl(\mathbf S(\mu_1,\mu_2),\mathbf A(\mu_1,\mu_2)\bigr).
\end{equation}
Only the moments with $n\in\mathcal O$ enter $\Theta$; it is not asserted
to be an injective coordinate chart of probability measures.

For each $p\ge 0$, we define the active ambient moment polynomial. In these coordinates, the active fitting residual becomes an ambient polynomial
in \((\mathbf S,\mathbf A)\).
\begin{equation}
\label{eq:ambient-active-reduced-M}
    \widetilde{\mathcal M}_p(\mathbf S,\mathbf A)
    \coloneq
    \sum_{\substack{n,m\in \mathcal O\\ 0\le r\le \min(n,m)\\ p=n+m-2r}}
    c_{n,m,r}(a_n a_m^\top)\otimes
    \sym{(S_n+A_n)\contr{r}(S_m-A_m)}.
\end{equation}
The corresponding ambient active fitting functional is
\begin{equation}
\label{eq:ambient-active-fit-functional}
    \widetilde{\mathcal J}_{\mathrm{fit}}(\mathbf S,\mathbf A)
    \coloneq
    \frac12\sum_{p\ge 0}
    \big\|
        \widetilde{\mathcal M}_p(\mathbf S,\mathbf A)-\mathcal M_p^o
    \big\|_{\HS}^2.
\end{equation}
\end{definition}

Membership in either branch implies that its finite active moments lie
on the corresponding coordinate axis.
\begin{definition}
For any measurable set $A$, let $-A \coloneq \{-w \mid w\in A\}$.
Consider the following two branches:
\begin{enumerate}
\item \textbf{Antisymmetric branch}: The antisymmetric branch is the candidate path toward the antipodal saddle.
\begin{equation}
    \submanifoldforsaddlepoint
    \coloneq
    \left\{
        \mu_1\otimes \mu_2
        \mid
        \mu_1(-A)=\mu_2(A)
        \text{ for every measurable }A
    \right\},
\end{equation}
which is equivalent to
\begin{equation}
    T_{1,n}=(-1)^nT_{2,n}
    \qquad
    (n\ge 0)
\end{equation}
and this implies $\mathbf S=0$.
\item \textbf{Symmetric branch}: The symmetric branch contains the same-density teacher configuration.
\begin{equation}
    \mathcal M_+
    \coloneq
    \left\{
        \mu_1\otimes \mu_2
        \mid
        \mu_1=\mu_2
    \right\},
\end{equation}
which is equivalent to
\begin{equation}
    T_{1,n}=T_{2,n}
    \qquad
    (n\ge 0)
\end{equation}
and this implies $\mathbf A=0$.
\end{enumerate}

We write the fitting loss restricted to each branch as a one-coordinate functional.
\begin{equation}
\label{eq:branch-reduced-functional-Mplus-block-v3}
    \mathcal J_+(\mathbf S)
    \coloneq
    \widetilde{\mathcal J}_{\mathrm{fit}}(\mathbf S,0),
    \quad 
    \mathcal J_-(\mathbf A)
    \coloneq
    \widetilde{\mathcal J}_{\mathrm{fit}}(0,\mathbf A).
\end{equation}
\end{definition}

The equivalences above use all raw orders $n\ge0$, not only
$n\in\mathcal O$. Here is the justification on the compact sphere.
Equality of all tensor moments gives equality of every monomial integral
and hence every polynomial integral. Restricted polynomials form an
algebra containing constants and separating points of the sphere.
By Stone--Weierstrass, for each continuous $f$ and $\epsilon>0$
there is a polynomial $P$ with $\sup_\Sphere|f-P|<\epsilon$. Thus,
for two measures with the same raw moments,
\[
 \left|\int f\,d\mu_1-\int f\,d\mu_2\right|
 \le\int|f-P|\,d\mu_1+\int|f-P|\,d\mu_2<2\epsilon.
\]
Letting $\epsilon\downarrow0$ gives equality of the measures.
For reflection $\mathsf R(w)=-w$, the identity
\[
 T_n(\mathsf R_\#\mu_1)
 =\int(-w)^{\otimes n}\,d\mu_1(w)=(-1)^nT_n(\mu_1)
\]
gives the antisymmetric equivalence by the same argument.
Consequently $\mu\in\mathcal M_-$ implies $\mathbf S=0$ and
$\mu\in\mathcal M_+$ implies $\mathbf A=0$, but neither converse is
inferred from the finite active blocks.

The first proposition identifies the distinguished active-origin point where the two branches meet.
\begin{proposition}[Active origin fiber and the distinguished Dirac saddle]
\label{prop:intersection-zero-output-fiber}
\textbf{Assumptions used.} \cref{assumption:on-sphere,assumption:odd-symmetric-teacher}.

Define the active-coordinate zero fiber
\begin{equation}
    \mathcal C_0
    \coloneq
    \left\{
        \mu_1\otimes \mu_2
        \mid
        \Theta(\mu_1,\mu_2)=(0,0)
    \right\}.
\end{equation}
Then
\begin{equation}
    \mu_{\dag}
    \coloneq
    \nu_\star^\pm \otimes \nu_\star^\pm
    \in
    \mathcal M_+\cap\submanifoldforsaddlepoint\cap\mathcal C_0.
\end{equation}
In addition, if
\begin{equation}
    \mu_1=\mu_2=\nu_\star^\pm,
\end{equation}
then the Wasserstein gradients vanish. In other words,
$\mu_{\dag}=\nu_\star^\pm\otimes\nu_\star^\pm$ is a first-order stationary point.
\end{proposition}

\begin{proof}
The proof only uses antipodal symmetry, oddness of the active moments, and the flatness of the metric regularizer at the antipodal target.

$\mu_\dag$ is antipodally symmetric, equal in both factors, and has zero odd moments; hence it belongs to the displayed intersection.

If $\mu\in\mathcal C_0$, then $T_{1,n}=T_{2,n}=0$ for all $n\in\mathcal O$. Therefore, on $\mathcal C_0$, the fitting first variation vanishes because of the definition of $G_{i,n}^{\mathrm{fit}}$.

The full spherical gradient is only the regularization contribution
\begin{equation}
    \nabla_{\Sphere}\fv{F}{\mu_i}
    =
    \lambdaReg Z_{\pm}^{p_W/2-1} \nabla_{\Sphere} c_\star (w),
    \qquad i=1,2.
\end{equation}
If $\mu_1$ and $\mu_2$ are $\nu_\star^{\pm}$, then $Z_{\pm}=0$ and the spherical gradient vanishes.
\end{proof}

The next proposition fixes the opposite endpoint of the cartoon: the exact-fit teacher lies on the same-density branch.
\begin{proposition}[The reference exact-fit point lies on the same-density branch]
\label{prop:reference-fit-minimum-on-Mplus}
\textbf{Assumptions used.} \cref{assumption:on-sphere,assumption:odd-symmetric-teacher}.

The fitting part satisfies
\begin{equation}
    \widetilde{\mathcal J}_{\mathrm{fit}}(\mathbf T^o,0)=0
    =
    \inf_{(\mathbf S,\mathbf A)\in \mathcal H_{\mathcal O}\times \mathcal H_{\mathcal O}}
    \widetilde{\mathcal J}_{\mathrm{fit}}(\mathbf S,\mathbf A).
\end{equation}
Moreover, by the symmetric-teacher assumption, the reference teacher
product measure belongs to the same-density branch:
\begin{equation}
    \mu^o=\mu_1^o\otimes\mu_2^o\in \mathcal M_+,
    \qquad
    \mathbf S(\mu^o)=\mathbf T^o,
    \qquad
    \mathbf A(\mu^o)=0.
\end{equation}
In particular, the distinguished zero-fit global minimum of the data term lies on $\mathcal M_+$.
\end{proposition}

\begin{proof}
By \cref{assumption:odd-symmetric-teacher}, the teacher is symmetric across the two factors, so $\mu_1^o=\mu_2^o$ and therefore $\mu^o\in \mathcal M_+$. In the active coordinates this is exactly the identity $\mathbf S(\mu^o)=\mathbf T^o$ and $\mathbf A(\mu^o)=0$.

Next, by the definition of $\mathcal M_p^o$ and of the ambient moment map,
\begin{equation}
    \widetilde{\mathcal M}_p(\mathbf T^o,0)=\mathcal M_p^o
    \qquad
    \forall p\ge 0.
\end{equation}
Hence $\widetilde{\mathcal J}_{\mathrm{fit}}(\mathbf T^o,0)=0$. Since $\widetilde{\mathcal J}_{\mathrm{fit}}$ is a sum of squared Hilbert norms, it is nonnegative everywhere.
\end{proof}
